Quatre transformations, une seule question : que conservent-elles ? La translation et les
symétries conservent les distances, donc les angles et les aires ; l'homothétie de rapport
$k$ multiplie les distances par $|k|$ et conserve les angles.

Les démonstrations sont courtes parce que les transformations sont définies par leur
formule et non par une construction : la translation de vecteur $v$ envoie $x$ sur $x + v$,
et il suffit de calculer la différence de deux images. C'est exactement ce qu'un élève fait
avec des coordonnées, sans coordonnées.

Ce chapitre est aussi celui où l'écart entre l'énoncé scolaire et l'énoncé formel est le
plus visible : \og{} construire l'image d'une figure par une rotation \fg{} est une tâche,
pas une proposition. Ce qui se démontre, ce sont les propriétés de l'application ; le tracé
reste au tableau.
